\subsection{Evaluation Methodology}\label{sec:evaluation}
$G$ was evaluated on the test split. Each row of the window matrix consists of a condition window of $l$ realized returns and the realized return of the next day. The conditions were passed to the trained $G$ and the generated values were compared to the realized values. For each condition, $B=1000$ samples were drawn from $G$ with different noise vectors. 

The realized returns of the split, $x = (x_1, \ldots, x_{n_{\mathrm{cond}}})$, give the reference values, where $n_{\mathrm{cond}}$ is the number of rows in the window matrix. Since the first $l$ returns are used as the initial condition window, a split of $n$ returns gives $n_{\mathrm{cond}} = n - l$ rows.
The generated draws form a matrix $\hat{X} \in \mathbb{R}^{n_{\mathrm{cond}} \times B}$, where $\hat{x}_{i,j}$ is the $j$-th draw of condition $i$. Additional sets were derived from $\hat{X}$. 

The path set $\hat{x}_{\mathrm{path}} = (\hat{x}_{1,1}, \ldots, \hat{x}_{n_{\mathrm{cond}},1})$ is the first column of $\hat{X}$, giving one generated return per trading day.
The pooled set takes all $n_{\mathrm{cond}}B$ generated entries of $\hat{X}$.

\begin{equation}
    \hat{x}_{\mathrm{pool}} = \{\, \hat{x}_{i,j} \;|\;
    i = 1, \ldots, n_{\mathrm{cond}}, \; j = 1, \ldots, B \,\}.
\end{equation}
 
The mean set holds the mean of all draws for each condition.
\begin{equation}
    m = (m_1, \ldots, m_{n_{\mathrm{cond}}}), 
    \text{ where }
    m_i = \frac{1}{B} \sum_{j=1}^{B} \hat{x}_{i,j}.
\end{equation}
$\hat{x}_{\mathrm{pool}}$ contains $B$ times more observations than $x$, therefore giving a more reliable estimate of the tails. $\hat{x}_{\mathrm{path}}$ preserves the order in time, since it holds one generated return per trading day.

Distributional statistics were computed for $x$ and $\hat{x}_{\mathrm{pool}}$: mean, standard deviation, skewness, excess kurtosis, minimum, maximum, \SI{1}{\percent} and \SI{99}{\percent} quantiles. Skewness and excess kurtosis measure the asymmetry and the tail weight described in Sec.~\ref{sec:stylized-facts}. Additionally, the fractions of returns exceeding \SI{5}{\percent} and \SI{10}{\percent} in absolute value were computed. The fraction of observations whose absolute value exceeds $u$ is denoted as $P(|x| > u)$. The exceedance fractions measure the frequency of extreme returns, which determines the value of out-of-the-money options (Sec.~\ref{sec:options}).

$C_2(\tau)$ from Eq.~\ref{eq:vol-clustering} was computed for $x$ and $\hat{x}_{\mathrm{path}}$ at $\tau=1, 2, 5, 10$. $\hat{x}_{\mathrm{pool}}$ cannot be used since it has no order in time.

Lastly, the mode collapse described in Sec.~\ref{sec:fin-gan} was checked. For distributional mode collapse, the standard deviation of each row of $\hat{X}$ was computed. The smallest value was compared with $\epsilon=0.0002$. For mean mode collapse, the standard deviation of $m$ was compared with the same $\epsilon$ value. The evaluation is implemented in \texttt{ttf\_eval.py}.